% ---------------------------------------------------------------------------
% Chapter skeleton — matches chapters/01-shape-of-a-spark-application.tex.
% Replace every <...>. Delete these comments. Draft figures BEFORE prose.
% Targets per chapter: 6–9k words · ≥3 figures · ≥4 listings · ≤2 callouts/1k words
% ---------------------------------------------------------------------------
\chapter{<Title Case Chapter Title>}
\label{ch:<short>}

\begin{chapterintro}
<The specimen, as the reader meets it in this chapter's terms: a symptom with a
number. One sentence on what this chapter will NOT do. One sentence on the
payoff: what the reader can see or change by the end.>
\end{chapterintro}

\section{<The mechanism, in plain words>}
\label{sec:<short>-<topic>}

% Beat 1 — symptom on the specimen / in the UI
<Two or three sentences. A number from the reference cluster.>

% Beat 2 — the picture, introduced before it appears
<What the figure shows and where to look> (\Cref{fig:<slug>}).

\begin{figure}[htbp]
  \centering
  \includegraphics[width=\linewidth]{<slug>.pdf}
  \caption{<Takeaway sentence.> <How to read it: what red / dashed / green mean here.>}
  \label{fig:<slug>}
\end{figure}

% Beat 3 — name it
A \emph{<term>} is <one-sentence definition>.

% Beat 4 — prove it: listing + real output + arithmetic
\begin{lstlisting}[language=SQL,caption={<What this listing demonstrates>},captionpos=b,label=lst:<short>-<topic>]
EXPLAIN FORMATTED
SELECT ...
FROM orders o
...;
\end{lstlisting}

\begin{lstlisting}[language=text,numbers=none,caption={<Its output, trimmed>},captionpos=b,label=lst:<short>-<topic>-out]
== Physical Plan ==
AdaptiveSparkPlan (12)
+- ...
-- ... 3 Filter nodes elided
\end{lstlisting}
% CAPTURE: run on the reference cluster (App. A) and paste the real output.

Put numbers on it: <arithmetic on the reference cluster — 40 task slots, 16 GB executors>.

% Beat 5 — the lever and its cost
<Which config / rewrite / layout change moves the cost, with \conf{key} and its
3.5.2 default, and what the change costs in turn.>

\begin{tip}
<One diagnostic shortcut — a UI panel, a metric, a one-liner.>
\end{tip}

\section{<Next mechanism>}
\label{sec:<short>-<topic2>}
<Same five beats.>

\section{Summary}
\label{sec:<short>-summary}

<One prose paragraph restating the chapter's mental model in the terms the
reader can now use — every term from the chapter, no new ones.>

<One or two sentences: what the specimen still does wrong, and which chapter fixes it.>

\exercisehead
\begin{exercises}
\item <Run-and-observe on orders: do X, open the Stages tab, record Y.>
\item <Read-this-plan: given an EXPLAIN excerpt, identify Z.>
\item <Predict-then-check: predict the effect of changing \conf{...}, then verify.>
\end{exercises}
